% !TEX program = pdflatex
% =====================================================================
% modelo.tex -- esqueleto de documento avulso da linha Solara
%
% Copie este arquivo, renomeie, troque os tres campos da `\abertura` e
% escreva. Nao ha decisao de forma a tomar: mancha, fontes,
% paleta e abertura saem do `estilo-reporte.sty` e do `config/preambulo.tex`.
%
%   cp modelo.tex meu-registro.tex
%   latexmk -pdf meu-registro.tex
%
% MOTOR. pdflatex. A fonte e a Charis SIL, com matematica do newtxmath e
% Bera Mono, todas do TeX Live; o preambulo base recusa outro motor. O
% `% !TEX program` no topo faz o VS Code acertar o motor sozinho.
%
% DIRETORIO. Compile de dentro de `report/`. O caminho do preambulo base
% e o da marca sao relativos a ele.
%
% FORMA. E a de manual tecnico: coluna unica de 152 mm, titulos recuados
% 14 mm para a esquerda, rodape com o nome do arquivo, a data e "pagina X
% de Y". Avisos (`perigo`, `atencao`, `cuidado`, `aviso`), `nota`,
% `passos` e `historico` vem do preambulo base.
%
% CLASSE. E `book` porque a folha de estilo usa `\chaptertitlename` e os
% comandos de capitulo do `tocloft`; em `article` ela nao carrega.
% Documento avulso nao abre capitulo -- use `\section*` e
% `\subsection*`, como abaixo.
%
% DEITADO. Documento que existe para caber uma tabela larga corre dentro
% de `landscape` -- `tarefas-dados.tex` e `modelos-metodos.tex` sao os
% dois casos. La, carregue `geometry` proprio antes da folha de estilo e
% ponha `\pagestyle{empty}`, porque o rodape sai deitado na margem
% lateral depois que o `pdflscape` gira a folha.
% =====================================================================
\documentclass[10pt,a4paper,oneside]{book}

\usepackage[brazilian]{babel}
\usepackage{estilo-reporte}

% A pagina termina onde o texto termina.
\raggedbottom

\autoria{João Leonardi da Silva Melo}{joao\_leonardi.melo@somosicev.com}
% \marca{outro/arquivo.png}   % descomente so para trocar a marca

\begin{document}

\abertura{Título do documento}
         {Registro de execução · \today{} ·
          \texttt{caminho/do/script.py}}

\section*{Primeira seção}

Texto. O corpo e a Charis SIL a 10 pt, e a linha de 152 mm corre a 98,1
caracteres.

\campo{Rótulo} Parágrafo com rótulo destacado na cor da marca.

\termo{Termo} marca vocabulário novo; \exper{nome-do-experimento} marca
nome de experimento; \texttt{arquivo.py} marca caminho e nome de coluna.

\begin{ficha}
  \item[Primeiro]  a lista de ficha e para definição curta, e o rótulo
                   fica na sem serifa em negrito.
  \item[Segundo]   o corpo dela cai para 7,4 pt.
\end{ficha}

\pepar[Pendente]{a marca de vaga imprime o que falta e some quando o
campo e preenchido; uma varredura por \texttt{pe} diz quanto do
documento existe.}

\section*{Resultado}

\tabfont
\begin{tabular}{@{}>{\rag\arraybackslash}p{3.3cm}rr@{}}
\toprule
\cab{Modelo} & \cab{Viés} & \cab{MAE} \\
\midrule
Uma linha      & $-1{,}032$ & 1,301 \\
Outra linha    & $+0{,}258$ & 0,962 \\
\bottomrule
\end{tabular}
\notas{A nota de tabela cai para 6,8 pt e diz unidade, recorte e
janela. \num{271} usinas, 2025.}

\end{document}
